(1) $\sigma^2 = (v^*v)^2 + (vv^*)^2 - v^*v\,vv^* - vv^*\,v^*v = v^*v + vv^* - 0 - 0 = 1$, using idempotency (Proposition \claimref{singularity.descent-what-it-erases}) and $v^2 = (v^*)^2 = 0$ to remove the cross terms. In the atom (\claimref{singularity.pseudo-inverse-identities-hyperbolic-atom}): $v v^* = \tfrac14(e_0 + e_1)(e_0 - e_1) = \tfrac14(1 - e_0e_1 + e_1e_0 + 1) = \tfrac12(1 - e_0e_1)$, so $v^*v = 1 - vv^* = \tfrac12(1 + e_0e_1)$ and $\sigma = e_0 e_1$. (2) From $\sigma v = -v$ and $v\sigma = v$ (in the atom: $e_0e_1(e_0 + e_1) = -(e_1 + e_0)$ and $(e_0+e_1)e_0e_1 = e_1 + e_0$; abstractly from \eqref{eq:car}), and $\sigma^2 = 1$ giving $e^{-\lambda\sigma} = \cosh\lambda - \sigma\sinh\lambda$ (as for any element of square one; in the atom $\sigma = e_0 e_1$ is a bivector with scalar square, the case of \claimref{std-algebra.bivector-exponential}):
\[
e^{\lambda\sigma/2}\,v\,e^{-\lambda\sigma/2} = v\,e^{-\lambda\sigma} = v\cosh\lambda - v\sigma\,\sinh\lambda = (\cosh\lambda - \sinh\lambda)\,v = e^{-\lambda}v ,
\]
using $v\sigma = v$ once; mirror-symmetrically for $v^*$. (3) is $v^*v = \tfrac12(1 + \sigma)$. (4) $N = v^*v$ counts occupation in the two-state module, and $\sigma = 2v^*v - 1$. (5) The even subalgebra of $\Cl(1,1)$ is $\operatorname{span}\{1, \sigma\}$ with $\tilde\sigma = -\sigma$, so the rotors are the elements $a + b\sigma$ with $(a + b\sigma)(a - b\sigma) = a^2 - b^2 = 1$: a one-dimensional group whose tangent line at $1$ is $\R\sigma$. A continuous one-parameter group of rotors is therefore $\lambda \mapsto \exp(c\lambda\sigma/2)$ for a real $c$, which by the closed form (\claimref{std-algebra.bivector-exponential}) is $\cosh\tfrac{c\lambda}{2} + \sigma\sinh\tfrac{c\lambda}{2}$: the boost flow at a rescaled parameter. The rotor group itself has two components, $a \ge 1$ and $a \le -1$; the exponentials fill only the first, and $-1$ is a rotor that is no exponential. A one-parameter flow stays in the identity component, so the uniqueness of (5) is unaffected; it is not the statement that every rotor of the atom is $\exp(\lambda\sigma/2)$.
